% !TEX root = ../../../main.tex
\section{Complex eigenvalues}
\label{sec:linear-algebra-i:complex-eigenvalues}

% Sources: IMG_9326--IMG_9327; October 15 PDF, pp. 17--21.
\subsection{Complex arithmetic}

A complex number has the form \(z=x+iy\), with \(x,y\in\R\) and
\(i^2=-1\). Its real and imaginary parts are \(x\) and \(y\), respectively.
\begin{marginfigure}
  \centering
  % Generic complex number and its conjugate; the positive quadrant is illustrative.
\begin{tikzpicture}[
  x=10mm,y=10mm,font=\footnotesize,line width=0.45pt,
  axis/.style={black!40,line width=0.35pt,-{Latex[length=1.6mm,width=1.1mm]}},
  guide/.style={dashed,black!45,line width=0.35pt}
]
  \coordinate (Z) at (1.7,1.3);
  \coordinate (C) at (1.7,-1.3);
  \draw[axis] (-0.4,0) -- (2.8,0) node[right=2pt] {\(\operatorname{Re}\)};
  \draw[axis] (0,-1.8) -- (0,1.9) node[above=2pt] {\(\operatorname{Im}\)};
  \draw[guide] (0,1.3) -- (Z) -- (C) -- (0,-1.3);
  \draw (0,0) -- (Z) node[midway,above left=1pt] {\(\lvert z\rvert\)};
  \draw[black!55] (0,0) -- (C);
  \fill (Z) circle[radius=1.5pt];
  \fill (C) circle[radius=1.5pt];
  \node[above right=2pt] at (Z) {\(z=x+iy\)};
  \node[below right=2pt] at (C) {\(\overline z=x-iy\)};
  \node[left=3pt] at (0,1.3) {\(y\)};
  \node[left=3pt] at (0,-1.3) {\(-y\)};
  \node[above right=1pt] at (1.7,0) {\(x\)};
  \node[below left=2pt] at (0,0) {\(0\)};
  \path (current bounding box.south) ++(0,-2mm);
\end{tikzpicture}

  \caption[Conjugation in the complex plane]{A schematic for \(x,y>0\).
    Conjugation reflects a point across the real axis. The distance from
    zero is \(\lvert z\rvert=\sqrt{x^2+y^2}\) for both points.}
  \label{fig:linear-algebra-i:complex-conjugation}
\end{marginfigure}
Powers of \(i\) repeat with period four,
\[
  i^2=-1,\qquad i^3=-i,\qquad i^4=1,\qquad i^5=i,\qquad i^6=-1.
\]
The conjugate of \(z\) is \(\overline z=x-iy\), and its modulus is
the nonnegative number satisfying
\[
  |z|^2=z\overline z=x^2+y^2.
\]
For \(z\ne0\),
\[
  \frac1z=\frac{\overline z}{|z|^2}=\frac{x-iy}{x^2+y^2}.
\]
More generally, division by \(w\ne0\) can be written as
\(z/w=z\overline w/|w|^2\).

\autoref{fig:linear-algebra-i:complex-conjugation} identifies \(z\)
with the point \((x,y)\) and its conjugate with \((x,-y)\).

\subsection{Returning to the rotation example}

Regard the real matrix \(B\) from
\autoref{ex:linear-algebra-i:rotation-no-real-eigenvalue} as a matrix
acting on \(\C^2\). For \(\lambda=i\), its kernel equations give
\[
  v^1-iv^2=0\quad\Longleftrightarrow\quad v^2=-iv^1.
\]
For \(\lambda=-i\), they give \(v^1+iv^2=0\). Consequently,
\[
  \begin{aligned}
    E_i(B)&=\left\{\begin{bmatrix}t\\-it\end{bmatrix}\;\middle|\;t\in\C\right\},\\
    E_{-i}(B)&=\left\{\begin{bmatrix}t\\it\end{bmatrix}\;\middle|\;t\in\C\right\}.
  \end{aligned}
\]
In each case, the eigenvectors are obtained by taking \(t\ne0\).
% IMG_9327 has an isolated +1 in B-iI. The original matrix, its stated
% action, both kernel equations, and PDF p.18 fix that entry as -1.

% !TEX root = ../../hw06.tex
% Source: supplied Homework 6 PDF, page 1, Question 5.
\begin{exercise}
  \label{ex:hw06:rotation-eigenvalues}
  For any \(\alpha\in\R\), the real matrix
  \[
    R_\alpha=\begin{bmatrix}
      \cos\alpha&-\sin\alpha\\
      \sin\alpha&\cos\alpha
    \end{bmatrix}
  \]
  is called a rotation matrix, because it rotates each vector counterclockwise
  by the angle \(\alpha\). Compute the (possibly complex) eigenvalues and
  eigenvectors of \(R_\alpha\)
  \solutionlink{hw06:rotation-eigenvalues}
\end{exercise}


\subsection{Roots and multiplicities}

\begin{theorem}[Fundamental theorem of algebra]
  Every complex polynomial of degree \(n\ge1\),
  \[
    p(z)=a_nz^n+\cdots+a_0,\qquad a_n\ne0,
  \]
  factors as
  \[
    p(z)=a_n(z-\lambda_1)\cdots(z-\lambda_n),
  \]
  where \(\lambda_1,\ldots,\lambda_n\in\C\), with repetitions allowed.
\end{theorem}
% PDF p.19 states this theorem without a proof.

The characteristic polynomial of an \(n\)-by-\(n\) complex matrix
therefore has \(n\) roots counted with multiplicity. The multiplicity
of \(\lambda\) is the number of factors equal to \((z-\lambda)\).
The \(n\) eigenvalues need not be distinct. Since the convention here is
\(p_M(z)=\det(M-zI)\), its factorization has leading coefficient
\((-1)^n\).

\begin{example}
  For the identity matrix on \(\mathbb F^2\),
  \[
    p_I(x)=(1-x)^2,\qquad E_1(I)=\mathbb F^2.
  \]
  Its only eigenvalue is \(1\), with multiplicity two, and every
  nonzero vector is an eigenvector.

  In contrast, let
  \[
    J=\begin{bmatrix}1&1\\0&1\end{bmatrix}.
  \]
  It has the same characteristic polynomial, but
  \[
    (J-I)\begin{bmatrix}v^1\\v^2\end{bmatrix}
      =\begin{bmatrix}v^2\\0\end{bmatrix}.
  \]
  Hence
  \[
    E_1(J)=\left\{\begin{bmatrix}t\\0\end{bmatrix}\;\middle|\;t\in\mathbb F\right\},
  \]
  which has dimension one.
\end{example}
% PDF p.21 prints (0,c) after deriving v2=0. Use (c,0), as its
% displayed multiplication and the definition of the kernel require.

The different actions of \(I\) and \(J\) over \(\R\) are pictured in
\autoref{fig:linear-algebra-i:identity-and-shear-eigenspaces}.
\begin{figure}[htbp]
  \centering
  % The supplied matrices I and J=[1 1;0 1]. Equal units in both panels.
\begin{tikzpicture}[
  x=13mm,y=13mm,font=\footnotesize,line width=0.45pt,
  vector/.style={-{Latex[length=1.6mm,width=1.1mm]}},
  axis/.style={vector,black!40,line width=0.35pt},
  guide/.style={dashed,black!45,line width=0.35pt}
]
  \begin{scope}
    \node at (1.1,2) {\(I\)};
    \fill[black!8] (0,0) -- (1,0) -- (1,1) -- (0,1) -- cycle;
    \draw[axis] (-0.4,0) -- (2.6,0) node[right=2pt] {\(x\)};
    \draw[axis] (0,-0.3) -- (0,1.6) node[above=2pt] {\(y\)};
    \draw (1,0) -- (1,1) -- (0,1);
    \draw[vector] (0,0) -- (1,0) node[below=3pt] {\(e_1\)};
    \draw[vector] (0,0) -- (0,1) node[left=3pt] {\(e_2\)};
    \node at (1.1,-0.9) {\(E_1(I)=\R^2\)};
  \end{scope}
  \begin{scope}[xshift=61mm]
    \node at (1.1,2) {\(J\)};
    \fill[black!8] (0,0) -- (1,0) -- (2,1) -- (1,1) -- cycle;
    \draw[axis] (-0.4,0) -- (2.6,0) node[right=2pt] {\(x\)};
    \draw[axis] (0,-0.3) -- (0,1.6) node[above=2pt] {\(y\)};
    \draw[guide] (0,0) -- (0,1) -- (1,1) -- (1,0);
    \draw (1,0) -- (2,1) -- (1,1);
    \draw[vector] (0,0) -- (1,0) node[below=3pt] {\(Je_1=e_1\)};
    \draw[vector] (0,0) -- (1,1) node[above left=3pt] {\(Je_2\)};
    \node at (1.1,-0.9) {\(E_1(J)=\operatorname{span}\{e_1\}\)};
  \end{scope}
\end{tikzpicture}

  \caption[Equal multiplicities, different eigenspaces]{The identity
    fixes every vector. The shear \(J\) fixes only the horizontal axis,
    although both characteristic polynomials are \((1-x)^2\).
    The dashed unit square on the right shows the shape before applying
    \(J\); the shaded parallelogram is its image.}
  \label{fig:linear-algebra-i:identity-and-shear-eigenspaces}
\end{figure}

\begin{proposition}
  If \(\lambda\) is an eigenvalue of a complex matrix and has
  multiplicity \(m_\lambda\) in its characteristic polynomial, then
  \[
    1\le\dim E_\lambda\le m_\lambda.
  \]
\end{proposition}
% PDF p.21 states the upper bound without supplying a proof.

% !TEX root = ../../hw06.tex
% Source: supplied Homework 6 PDF, page 1, Question 3, parts 3.1--3.4.
% Preserve the unrestricted field in 3.2 and the absent field in 3.1;
% the truth values are for the reader to assess, not assumptions to repair.
\begin{exercise}
  \label{ex:hw06:eigenvalue-true-or-false}
  True or false? Explain. If false, provide a counterexample.
  \begin{enumerate}[{3}.1.][4]
    \item If a matrix has one eigenvector, then it has infinitely many
      eigenvectors.
    \item Every matrix over a field \(\mathbb F\) has at least one eigenvalue
      in \(\mathbb F\).
    \item Similar matrices always have the same eigenvectors.
    \item If two matrices have the same eigenvalues, then they are similar
  \end{enumerate}
  \solutionlink{hw06:eigenvalue-true-or-false}
\end{exercise}

